
\documentclass[10pt]{article}
\usepackage[utf8]{inputenc}
\usepackage[T1]{fontenc}
\usepackage{amsmath, amssymb, xcolor, geometry, enumitem, multicol, tcolorbox}

\definecolor{mainblue}{RGB}{0, 102, 204}
\geometry{a4paper, margin=1.2cm, top=1cm, bottom=3cm, footskip=1.2cm}

\begin{document}
\enlargethispage{2.5cm}

% Modern Header
\begin{tcolorbox}[colback=mainblue!5, colframe=mainblue, sharp corners, boxrule=0.5pt]
    \small \textbf{Unit:} undefined \hfill \textbf{Name:} \rule{3cm}{0.4pt} \\
    \small \textbf{Topic:} Variables and Expressions / Order of Operations \hfill \textbf{Date:} \rule{2cm}{0.4pt} \\
    \centering \textbf{\large Advanced (Level 1)}
\end{tcolorbox}

\vspace{0.2cm}

% MCQ Section
\begin{multicols}{2}
\begin{enumerate}[label=\textbf{Q\arabic*.}, leftmargin=*, itemsep=6pt, parsep=0pt]

  \item \small What is the variable in the expression $2x + 3$?
  \begin{enumerate}[label=(\Alph*), leftmargin=0.6cm, nosep, itemsep=2pt]
    \item 2
    \item x
    \item 3
    \item 2x
    \item + 
  \end{enumerate}

  \item \small Identify the coefficient in the term $4y$?
  \begin{enumerate}[label=(\Alph*), leftmargin=0.6cm, nosep, itemsep=2pt]
    \item y
    \item 4
    \item 4y
    \item 1
    \item None
  \end{enumerate}

  \item \small What operation is represented by the symbol $+$ in $2 + 5$?
  \begin{enumerate}[label=(\Alph*), leftmargin=0.6cm, nosep, itemsep=2pt]
    \item Subtraction
    \item Addition
    \item Multiplication
    \item Division
    \item Exponentiation
  \end{enumerate}

  \item \small In the expression $3(2x - 1)$, what is the term inside the parentheses?
  \begin{enumerate}[label=(\Alph*), leftmargin=0.6cm, nosep, itemsep=2pt]
    \item 3
    \item 2x
    \item 2x - 1
    \item 3(2x)
    \item None
  \end{enumerate}

  \item \small What is the constant term in $2x^2 + 3x - 4$?
  \begin{enumerate}[label=(\Alph*), leftmargin=0.6cm, nosep, itemsep=2pt]
    \item 2
    \item 3
    \item -4
    \item 2x
    \item 3x
  \end{enumerate}

  \item \small Identify the operator in $7 - 2$?
  \begin{enumerate}[label=(\Alph*), leftmargin=0.6cm, nosep, itemsep=2pt]
    \item +
    \item -
    \item *
    \item /
    \item ^
  \end{enumerate}

  \item \small What is the exponent in $x^3$?
  \begin{enumerate}[label=(\Alph*), leftmargin=0.6cm, nosep, itemsep=2pt]
    \item x
    \item 3
    \item x^3
    \item 1
    \item None
  \end{enumerate}

  \item \small In the expression $5(2x + 1)$, what is the factor?
  \begin{enumerate}[label=(\Alph*), leftmargin=0.6cm, nosep, itemsep=2pt]
    \item 2x + 1
    \item 5
    \item 5(2x)
    \item 2x
    \item None
  \end{enumerate}

\end{enumerate}
\end{multicols}

\vspace{-0.2cm}
\hrule
\vspace{0.2cm}
\textbf{\small Open Ended Questions (FRQ)}
\begin{enumerate}[label=\textbf{Q\arabic*.}, start=9, itemsep=5pt, topsep=0pt]

  \item \small Draw a diagram to illustrate the concept of the order of operations in the expression $3 + 4 \times 2$. Explain your reasoning. \vspace{2cm}

  \item \small Explain in words the difference between a variable and a constant in an algebraic expression. Provide an example of each. \vspace{2cm}

\end{enumerate}

\vfill

% Chef's Tips Footer
\begin{tcolorbox}[colback=blue!5, colframe=mainblue!50, title=\small \textbf{Chef's Tips}, fonttitle=\bfseries, bottom=1mm]
\begin{itemize}[noitemsep, topsep=2pt, leftmargin=0.5cm]
  \item \scriptsize Emphasize the importance of identifying and understanding the basic components of expressions, such as variables, constants, and operators.\item \scriptsize Use simple, one-step examples to reinforce the definitions and concepts, ensuring students grasp the foundational elements before progressing to more complex problems.\item \scriptsize Encourage students to visualize and explain their reasoning, especially for free-response questions, to ensure a deep understanding of the concepts rather than just memorization.
\end{itemize}
\end{tcolorbox}

\end{document}
  